% ---------------------------------------------------------------------------
% Methodology (Section 4) - drafted for M2. The evaluation metrics subsection
% is in evaluation_metrics.tex (M1). Citation keys match report/references.bib.
% Needs in the preamble: \usepackage{booktabs, amsmath}.
% All five approaches are final (A3-A5 configurations in configs/a{3,4,5}_best.yaml).
% ---------------------------------------------------------------------------
\section{Methodology}
\label{sec:method}

\subsection{Data and experimental protocol}
The challenge training set contains 4{,}200 labelled clips (350 per word), all 16~kHz mono 16-bit
PCM~\cite{zindi2021swahili}. The test labels are hidden, so we froze one stratified 70/15/15 split of the training
set (seed 42) into 2{,}940 training, 630 validation and 630 test clips (52--53 per word in validation and test). A
hash of the split is checked by every notebook. No byte-identical duplicate recordings were found.
Model selection uses the validation split only. The test split is scored once, after every configuration is frozen.
Experiments proceed in rounds: R0 establishes each approach, R1 trains default neural configurations, R2 tunes, R3
trains the final configurations with three seeds and evaluates them on the test split, and R4 varies the amount of
training data. Every run is logged with the change it makes and the reason for it (\texttt{results/experiments.csv}).

\subsection{Preprocessing}
Clips last 2.2--62~s (median 4.3~s), but a spoken digit lasts under a second, so most of each clip is silence or
background noise. A single silence-trimming threshold was unreliable: at 30~dB, 10\% of trimmed clips in a
600-clip sample still exceeded 4.5~s. We therefore keep the 1.5~s window with the highest signal energy. When the whole word fits, many
window positions hold almost the same energy, so we take the middle of that plateau. This centres the word: the
energy centroid lies within $\pm0.2$~s of the window centre for 90\% of clips. The window retains a median of
99.8\% of each clip's above-noise energy, and less than 70\% in only 2.7\% of clips, which are mostly long
recordings containing several utterances. Each window is peak-normalised.

\subsection{Acoustic features}
Frames use a 25~ms window and a 10~ms hop, giving 151 frames per window. We compute a 64-band log-mel spectrogram
(dB, 80~dB dynamic range) and mel-frequency cepstral coefficients (MFCCs)~\cite{davis1980comparison} with delta and
delta-delta coefficients (a 9-frame window). Features are computed once, cached, and shared by all members, so every
approach sees identical inputs. Normalisation statistics are always fitted on the training split only.

\subsection{Approaches}
Table~\ref{tab:approaches} summarises the five approaches. They were chosen to vary \emph{how} temporal order is
modelled: not at all (A1), by a generative Markov chain (A2), by recurrence (A3), by convolution over time (A4) and
by self-attention with self-supervised pretraining (A5).

\paragraph{A1: MFCC statistics with a linear classifier}
Each clip becomes one 156-dimensional vector: the mean, standard deviation, minimum and maximum over time of 13
MFCCs with their deltas. This discards temporal order entirely. After standardisation, a multinomial logistic
regression ($C=0.1$, chosen by 5-fold cross-validation on the training split) outperformed an RBF-kernel SVM with
Platt-scaled probabilities~\cite{platt1999probabilistic,wu2004probability}.

\paragraph{A2: GMM-HMM per word}
Following the classic isolated-word recogniser~\cite{rabiner1989tutorial}, each word has its own left-to-right HMM:
each state either repeats or moves to the next. Each state emits diagonal-covariance Gaussian mixtures over
39-dimensional frames (13 MFCCs with deltas). The frames are normalised per utterance (cepstral mean and variance
normalisation, CMVN) to remove the channel differences of about 300 speakers recording on different phones. States
are initialised by a flat start: every training sequence is split uniformly into segments, and state $i$ is
initialised from the $i$-th segments. Training uses Baum--Welch, with small MAP priors on the mixture weights,
variances and allowed transitions. These priors prevent components starved of frames from producing degenerate
(NaN) parameters. A clip is assigned to the word whose model gives the highest per-frame log-likelihood. The final
model has 12 states and 8 Gaussians per state.

\paragraph{A3: bidirectional LSTM with attention}
The log-mel frames ($151\times64$) feed two stacked bidirectional LSTM layers~\cite{hochreiter1997long,graves2013speech}
with 128 units per direction. The outputs are pooled over time with additive attention~\cite{bahdanau2015neural},
followed by dropout and a linear layer over the 12 words. Recurrence carries information across the whole word, and
attention learns which frames matter. It also exposes, for the error analysis, which frames each decision relied on.
The design follows the attention-based recogniser of de Andrade et al.~\cite{deandrade2018neural} with one
deliberate difference. We omit their convolutional layers by default, so that A3 is purely recurrent and contrasts
cleanly with the convolutional A4. Adding them back is one of the R2 ablations.

The R2 search selected:
\begin{itemize}
  \item 40 MFCCs with deltas instead of the log-mel input;
  \item a 1.0~s window;
  \item SpecAugment;
  \item a peak learning rate of $2\times10^{-3}$;
  \item de Andrade et al.'s convolutional front-end (two 1-D convolutions, kernel 5, 128 channels).
\end{itemize}
The final model has 855K parameters. The front-end lowered the validation log loss only marginally ($0.142\to0.134$)
and did not change the accuracy (97.0\%). We therefore also report the purely recurrent variant as the reference
point when comparing recurrence with convolution.

\paragraph{A4: temporal convolutional network (TC-ResNet)}
A4 follows TC-ResNet~\cite{choi2019temporal}. It treats the frequency axis as channels, and slides
one-dimensional convolutions over time only. A three-wide stem convolution feeds residual blocks. Each block has two
convolutions of kernel width 9, the first with stride 2, and a projected shortcut when the shape changes. Global
average pooling over time feeds a linear layer over the 12 words.

Our default is TC-ResNet8 at width 1.5: three blocks with 24--72 channels. It has 146K parameters, and a receptive
field of 171 frames (1.7~s), which already spans the whole input window. Temporal order therefore reaches A4 only
through the arrangement of its strided, stacked filters. There is no state carried through time, as there is in A3.
Because the classifier follows global average pooling, applying it at every time step gives an exact class
activation map over time, which we use for interpretation. The R2 search (below) kept this architecture, and
selected MFCC-40 with $\Delta$ and $\Delta\Delta$ (120 channels) from a 2.0~s window, with full augmentation. The
final model has 150K parameters (\texttt{configs/a4\_best.yaml}), 5.7 times fewer than the final A3.

\paragraph{A5: fine-tuned self-supervised speech transformer}
A5 fine-tunes a Transformer~\cite{vaswani2017attention} speech encoder that was pretrained without labels, following
wav2vec~2.0~\cite{baevski2020wav2vec}. A convolutional feature encoder turns the raw 16~kHz waveform of the same
centred window into one vector per 20~ms. A Transformer then contextualises these vectors, relating every frame to
every other through self-attention. Order enters only through a convolutional positional embedding. The frames of
the last layer are projected to 256 dimensions, mean-pooled over time, and classified.

The default backbone is XLS-R~\cite{babu2022xlsr} with 300M parameters, pretrained on 436K hours in 128 languages,
including 91 hours of Swahili. As is standard, the convolutional feature encoder stays frozen, and the Transformer
is fine-tuned. The new, randomly initialised head has its own, higher learning rate. The only built-in
regularisation is wav2vec~2.0's masking of latent frames: spans of 10 frames, started with probability 0.05. Both
backbones in the R2 comparison use these settings, although they were pretrained with different
defaults. The R2 search kept this setup and added waveform augmentation. The final model has 316M parameters, of
which 311.5M are fine-tuned (\texttt{configs/a5\_best.yaml}): 370 times more than A3 and 2{,}100 times more than A4.

\begin{table}[t]
  \centering
  \caption{The five approaches and how each treats temporal order.}
  \label{tab:approaches}
  \small
  \begin{tabular}{@{}llp{0.42\linewidth}@{}}
    \toprule
    ID & Model & Temporal order \\
    \midrule
    A1 & MFCC statistics + logistic regression & Discarded (bag of frames); $\Delta$ features keep $\sim$90~ms of local change \\
    A2 & GMM-HMM per word & Left-to-right Markov states \\
    A3 & BiLSTM + attention & Recurrence in both directions \\
    A4 & TC-ResNet & Convolution over time \\
    A5 & XLS-R (fine-tuned) & Self-attention; multilingual self-supervised pretraining \\
    \bottomrule
  \end{tabular}
\end{table}

\subsection{Training the neural models}
A3 and A4 share one training loop (\texttt{src/train\_torch.py}):
\begin{itemize}
  \item \emph{Optimisation:} AdamW with a one-cycle learning-rate schedule, gradient clipping at 1.0, cross-entropy
        with label smoothing 0.1, batch size 64, and at most 40 epochs.
  \item \emph{Early stopping:} training stops when the validation log loss (without label smoothing) has not
        improved for five epochs, and the weights of the best epoch are restored.
  \item \emph{Augmentation:} applied on the fly to the feature frames of training clips, using SpecAugment time and
        frequency masking~\cite{park2019specaugment}, time shifts of up to 100~ms and tempo perturbation
        (0.9--1.1$\times$). Noise injection at 10--30~dB SNR assumes a dB power scale, so it is used for log-mel
        inputs only. Padding and masks take the clip's minimum and mean value. For log-mel, whose bands share one
        dB scale, that is one value per clip. A4's MFCC inputs use one value per channel, because their channels
        differ in scale by orders of magnitude. With a single value, the first MFCC run of A4 collapsed to chance.
        All augmentation is synthetic, because the competition allows no external data.
  \item \emph{Reproducibility:} every generator is seeded before the weights are initialised, so a configuration
        reproduces exactly.
\end{itemize}
A5 uses the same loop, instead of a separate fine-tuning framework, so that the loss, model selection and
calibration are identical across the neural models. Four settings differ, as fine-tuning a pretrained model
requires:
\begin{itemize}
  \item each waveform is normalised to zero mean and unit variance inside the model, instead of standardised with
        training-split statistics;
  \item the pretrained Transformer is fine-tuned at a peak learning rate of $3\times10^{-5}$, and the new head at
        $10^{-3}$, with 10\% warm-up;
  \item training runs for 10 epochs with batch 16 and keeps the best epoch, with no early stopping;
  \item augmentation acts on the waveform: time shifts, speed perturbation (0.9--1.1$\times$, which unlike the
        log-mel tempo change also shifts the pitch), and noise at 10--30~dB SNR.
\end{itemize}
Every model's probabilities are then temperature-scaled~\cite{guo2017calibration}, as described in
Section~\ref{sec:metrics}.

\subsection{Hyperparameter search}
Classical models were tuned by grid search: cross-validated on the training split for A1, and on the validation
split for the HMM topology and preprocessing of A2. For A3 we ran a greedy, one-factor-at-a-time search in round R2.
Starting from the default configuration, each run changes a single factor relative to the current best
configuration, and keeps the change only if the validation log loss improves. The factors were augmentation,
pooling (attention, last state, mean), recurrent direction, input features (log-mel or MFCC), window length
(1.0/1.5/2.0~s), the way the word is located, the learning rate, and a convolutional front-end. A4 followed the
same protocol over the factors that apply to it: augmentation, input features, window length, the way the word is
located, depth, kernel width, network width and the learning rate. A5's search covered the pretraining language
(XLS-R against the English-only wav2vec~2.0 Base~\cite{baevski2020wav2vec}), freezing the whole encoder,
augmentation, the window length, a weighted sum of all layers instead of the last~\cite{yang2021superb}, and the
learning rate. This makes every decision traceable to a logged comparison. The cost is that interactions between factors are only explored along the chosen path. Runs that
failed, or that a later fix made stale, are kept outside the results table (\texttt{results/failed\_runs/} and
\texttt{results/superseded\_runs/}).

Single-seed comparisons are noisy, so we measured the noise. Retraining the final A3 configuration with three seeds
(42, 13, 7) gave a validation log loss of $0.137\pm0.004$ (mean $\pm$ SD) and an accuracy of $96.9\pm0.2\%$.
Training is deterministic: the seed-42 run reproduced the R2 result exactly. Search steps that changed the log loss
by about one SD are therefore within seed noise. These include the window length (1.0 versus 1.5~s), the learning
rate and the rejection of the 2.0~s window. We report them, but draw no conclusions from them. Only the larger
effects, such as augmentation, pooling, direction and the way the word is located, are interpreted.

A4 is much less stable across seeds: $0.135\pm0.023$ log loss and $96.8\pm0.3\%$ accuracy. One seed stopped
early, at epoch 16 of 40, while the learning rate was still high. For A4, therefore, only the effect of the way the
word is located (silence trim: $+0.157$ log loss) is clearly larger than seed noise. The other R2 differences
(0.02--0.04) are read as directions, not established effects. A5 is the most stable model: $0.1015\pm0.0009$ log
loss and $98.0\pm0.1\%$ accuracy. Pretrained weights leave the seed little to decide. So all of A5's R2 differences,
except the tie between the 1.5 and 2.0~s windows, are larger than its seed noise.

% TODO: add software references (PyTorch, scikit-learn, librosa, hmmlearn) after verifying them like the others.
